\documentclass[11pt]{article}
\usepackage[margin=1in]{geometry}
\usepackage{amsmath,amssymb,amsthm,booktabs,hyperref,xcolor,enumitem,graphicx}
\emergencystretch=3em
\hypersetup{colorlinks=true,linkcolor=blue,urlcolor=blue,citecolor=blue}
\newtheorem{proposition}{Proposition}
\newtheorem{remark}{Remark}
\title{A Certificate-Oriented Study of Covering the Unit Disk with 100 Equal Disks}
\author{Open computational research dossier\\\small springoxn/disk-covering-100}
\date{October 2026}
\begin{document}
\maketitle
\begin{abstract}
We study the minimum radius $r_D(100)$ for covering the closed unit disk by 100 congruent closed disks. The current certificate chain proves $r_D(100)>0.113235$ through six exhaustive boundary-count branches, while a separately replayed configuration gives $r_D(100)\le 0.1154469548855829$. The target $0.11324$ is not claimed: the $B=36$ branch retains a strict obstruction for the current pair relaxation. This dossier explains the geometric reduction, rational interval verification, upper and lower bounds, external dependencies, and unsuccessful routes, with each quantitative claim tied to a machine-readable artifact.
\end{abstract}
\section{Problem and scope}
Let $D=\{x\in\mathbb R^2:\|x\|\le 1\}$. We define
\[
r_D(100)=\inf\left\{r:\ D\subseteq\bigcup_{i=1}^{100}B(c_i,r),\quad c_i\in\mathbb R^2\right\}.
\]
The problem is continuous and non-convex. The centres are unrestricted, and the combinatorial type of the induced Voronoi or covering decomposition changes during optimization. Consequently, a numerical configuration supplies an upper bound, whereas a lower bound must exclude every admissible configuration at a prescribed radius.

This work is a certificate-oriented computational study rather than a claim of a formalized theorem in a proof assistant. The final decision layer uses integer, rational, or outward interval arithmetic whenever possible. Floating-point optimization and high-precision numerics are discovery tools only.

\subsection{What is proved and what is not}
The public release certifies
\[
  0.113235<r_D(100)\le 0.1154469548855829.
\]
The lower inequality is a six-branch exclusion certificate at $R=22647/200000$. The upper inequality is a replayed explicit configuration. Neither statement identifies the exact optimum. In particular, the nearby target $0.11324$ remains open in the present certificate system.

The distinction matters: a successful local box exclusion is not a global proof; an unblocked numerical point is not a certificate; and a covering configuration does not prove optimality.

\section{The upper-bound side}
An upper bound is constructive. One supplies 100 centres and proves that every point of $D$ lies within the proposed radius of at least one centre. The release contains two related forms of evidence.

\subsection{Replayed configuration}
The certificate `external-disk100-bounds-67ebf1c.json` records a replayed configuration with
\[
 r_D(100)\le \frac{1154469548855829}{10^{16}}.
\]
The verification decomposes the unit disk into boundary and interior cells, checks the relevant farthest points with exact or outward-rounded arithmetic, and binds the input and checker hashes. The result is an existence statement, not an optimality statement.

\subsection{A transparent lattice construction}
A separate certificate describes a triangular-lattice/hexagonal-shell construction with candidate radius $\alpha$ satisfying $\alpha^2=1/67$. Its independent checker verifies the retained lattice indices, shell counts, separation conditions, the positive root bracket, and irreducibility of $67t^2-1$. This construction is deliberately presented as a transparent upper-bound witness; it is not asserted to be the best 100-centre configuration.

\subsection{Why the upper bound is not the hard part}
The upper-bound workflow can optimize or replace the configuration and then re-run coverage checks. The lower-bound workflow cannot optimize over a finite list: it must cover arbitrary centre sets and all combinatorial types. This asymmetry is why the main effort is spent on exhaustive reductions and interval certificates.

\section{The six-branch lower-bound architecture}
Fix a candidate radius $R$ and assume, for contradiction, that 100 disks of radius $R$ cover $D$. The reduction classifies the boundary-serving count into
\[
 B31,\ B32,\ B33,\ B34,\ B35,\ B\ge 36.
\]
The first five branches reduce possible planar covering graphs to finite families using Euler relations, boundary ownership, triangulation, ears, chords, and internal-face counts $M$. Every accepted branch must provide (i) same-radius premises, (ii) exhaustive model coverage, (iii) an exact or outward-rational geometric inequality, and (iv) a positive budget margin.

\subsection{Finite combinatorial branches}
The B31--B35 certificates use different combinations of degree/curvature losses, core-ear exchange, ring models, chordless models, one-chord models, and special reductions for $M=25$ and $M=26$. The final evidence is stored in the B31, B32, B33, B34, and B35 JSON certificates. Their exclusion statements were established at a radius no smaller than the accepted candidate; radius monotonicity then transfers exclusion to smaller radii.

\subsection{The high-boundary branch}
For $B\ge36$, boundary arcs are arranged cyclically and adjacent pairs are bounded by a function of two half-arc parameters. A representative direct term has the form
\[
 \Phi(a,b)=\lambda\lVert c_a-c_b\rVert^2+\frac{H(a)+H(b)}2
 -\frac\mu2\bigl(\arcsin(s_a)+\arcsin(s_b)\bigr).
\]
The verifier partitions the parameter square into closed rational rectangles and proves an upper bound on every rectangle. The cyclic penalty uses the exact identity
\[
 \sum_i\frac{a_i+a_{i+1}}2=\sum_i a_i=\pi,
\]
not an unproved equal-arc Jensen assertion. The resulting area inequality has a B=36 margin and a positive increment per additional boundary cell, so the tail is closed once the B=36 pair domain is closed.

At $R=22647/200000$, the B36 replay visited 1,841,089 rational boxes and accepted 917,297; the three marked-side support classes used 3761, 104, and 35 boxes. These counts are certificate data, not numerical sampling claims.

\subsection{The target obstruction}
At $R=2831/25000=0.11324$, the current pair relaxation has an exact obstruction near $x=y=9921033/12500000$. Both support branches are zero there, so merely refining support intervals cannot repair the margin. The residual is therefore a geometric limitation of the present relaxation, not a missing grid refinement.

\section{Routes attempted: successes, failures, and open directions}
A useful research record must include routes that failed. They identify which relaxations are too weak and prevent rediscovering the same obstruction.

\subsection{Combinatorial and topological routes}
Dart-map, Voronoi-owner, triangulation, and generic-map compilers were used to turn geometric cases into finite combinatorial objects. Euler counting, ear decompositions, non-adjacent ear deletion, and core inequalities closed selected low-$M$ branches. Ring enumeration, chordless models, and one-chord models closed selected high-$M$ branches. The main failure mode was not enumeration itself but proving that every geometric degeneration was represented by the finite catalogue; this is why model ledgers and scope fields are retained.

\subsection{Capacity and local pruning}
Hall-style capacity trees, weighted witness capacities, axial multi-anchor boxes, spatial branch trees, reverse Voronoi distance bounds, annular cuts, radial hierarchies, directional counts, and shared-vertex controls all produced useful exclusions or ranked survivors. None alone closed the full target radius. Nonempty boxes remain unresolved; they are not silently counted as excluded.

\subsection{Energy, arc, and lens attempts}
Angular and harmonic energy proposals produced negative witnesses or counterexamples, showing that simple global energies were too weak. Arc caps, lens supports, Young-type support inequalities, and boundary-arc budgets became the successful B$\ge36$ route at $0.113235$. Replacing a coarse $4\lambda R^2$ distance cap by the exact centre distance was tested but is already contained in the original direct expression; it therefore adds no new geometric information. A common-endpoint lens inequality or a genuinely coupled multi-edge averaging inequality remains the principal open route for B36 at $0.11324$.

\subsection{Radius scanning and engineering}
A rational scan of 101 radii between $0.1132301$ and $0.11324$ diagnosed the collision between the old B36 budget and its exact obstruction. Candidates $0.113236$ and $0.113237$ were attempted with full pair subdivision but did not produce completed PASS or FAIL certificates before the runs were stopped. They must not be reported as mathematical failures. The next engineering improvement is a resumable queue of rational boxes with cached inverse-trigonometric and support intervals, independent hashes, and diagonal/non-diagonal sharding.

\section{Software, dependencies, and reproducibility}
The release uses Python's standard library for the final rational ledgers, especially `Fraction`, exact integer arithmetic, interval-style outward rounding, JSON manifests, and SHA-256 binding. SymPy is used for symbolic contracts and factor checks; NumPy and mpmath are used in discovery and high-precision exploration; optional python-flint/LLL tools support selected integer-relation experiments. Floating-point output is never sufficient for a final exclusion certificate.

The project also replays and independently audits material associated with VonEquinox's \texttt{DiskCoveringSolve}. That material is treated as an external reference and replay input, not as an automatically inherited proof. Third-party source is not vendored in this release.

The public package provides \texttt{reproduce.py}, which checks the core certificate files, their statuses, the candidate radius, the six-branch closure flag, and SHA-256 digests. The full route inventory is in \texttt{ROUTE\_INDEX.md} and \texttt{route\_manifest.json}; the claim-to-evidence table is in \texttt{docs/EVIDENCE\_LEDGER.md}.

A PASS is scoped: it certifies the assertions recorded in its JSON certificate, including its domain, radius, dependencies, and machine checks. It does not mean that every geometric reduction has been formalized in a proof assistant. This limitation is part of the result, not a footnote.

\section{Evidence ledger and conclusion}
\begin{table}[h]
\centering
\small
\resizebox{\linewidth}{!}{\begin{tabular}{p{0.28\linewidth}p{0.18\linewidth}p{0.43\linewidth}}
\toprule
Claim & Status & Primary artifact \\
\midrule
Explicit 100-disk configuration upper bound & replayed & \texttt{external-disk100-bounds-67ebf1c.json} \\
Transparent lattice candidate & independently checked & \texttt{independent\_lattice97\_upperbound\_n100.json} \\
B31--B35 exclusion & PASS & B31--B35 completion certificates \\
B$\ge36$ at $R=0.113235$ & PASS & \texttt{n100-tail-candidate-224.json} \\
Six-branch assembly at $R=0.113235$ & PASS & \texttt{n100-candidate-complete-225.json} \\
$R=0.11324$ under the old pair relaxation & unresolved/obstructed & support obstruction and direct-term diagnostics \\
$R=0.113236$ and $0.113237$ & unverified & stopped before a completed certificate \\
\bottomrule
\end{tabular}}
\caption{The status of the principal quantitative claims.}
\end{table}

The present rigorously supported interval is therefore
\[
 \boxed{0.113235<r_D(100)\le0.1154469548855829}.
\]
The lower endpoint is a certified exclusion radius; the upper endpoint is a certified explicit configuration radius. The gap is real mathematical work, not merely a formatting or precision issue. The public repository is designed to make that gap—and the routes that created it—inspectable.

\end{document}
